\chapter{Introduction}

\section{Motivation}
\ \ \ \ \, Federated learning (FL) trains a shared model across distributed clients
without requiring the clients to upload their raw data. This design is
attractive for edge intelligence because it reduces privacy exposure and
keeps data near its source device. A standard FL round includes broadcasting current global model, 
letting clients run local optimization and aggregates
their updates on a central server. The FedAvg algorithm is the classic
example of this workflow and remains the baseline for many FL systems
\citep{mcmahan2017communication}.

The communication cost of conventional digital FL grows with the number of
clients because each client must send an update to the server separately.
Analog over-the-air aggregation changes the communication model. If clients
transmit the updates in analog waveforms at the same time, the wireless
multiple-access channel physically superposes the signals, then server can
only observe an aggregate update in one shared channel use and fail to
decode each client separately. This makes OTA-FL a promising architecture
for large-scale wireless learning, yet with a main obstacle that 
same wireless channel can also injects noise, interference, and possible fading. 
The server therefore receives a corrupted
aggregate rather than the exact average update. When a plain FedAvg-style
server updates using this noisy vector directly, heavy-tailed impulsive
interference can cause slow convergence, oscillation, or divergence. The
reference work ``Adaptive Federated Learning Over the Air'' studies this
problem for adaptive methods under fading and alpha-stable interference
\citep{wang2024adaptive}. This project follows that direction by making the
alpha-stable tail index a core experimental variable; AWGN is treated as the
special case $\alpha=2$.

\section{Problem Statement}
\ \ \ \ \, This thesis studies FL over a simulated analog OTA aggregation channel.
At communication round $t$, the server receives a noisy aggregate update $\tilde{g}^t$:
\begin{equation}
    \tilde{g}^t =
    \frac{1}{N}\left(\sum_{n=1}^{N} \Delta_n^t + \xi^t\right),
\end{equation}
where $N$ denotes the number of clients, $\Delta_n^t$ denotes the update produced by
client $n$, and $\xi^t$ denotes symmetric alpha-stable channel noise.

When $\alpha<2$, smaller $\alpha$ means heavier tails and more frequent impulsive outliers, while $\alpha=2$ 
 recovers the Gaussian AWGN case. The central question is whether adaptive
server optimizers can use their historical first- and second-moment states to
make the training process more robust to this noisy aggregate.

The project should not be presented as a complete reproduction of
\citet{wang2024adaptive}. The original paper contains a full convergence
analysis and a broader experimental suite. The current codebase implements an
empirical simulator for the core heavy-tail mechanism: alpha-stable OTA
aggregation, adaptive server updates, and robust MAC pre-processing. Its
contribution is a controlled, reproducible software validation rather than a
new theoretical optimizer.

\section{Contributions}
The contributions of this project are as follows.
\begin{itemize}
    \item A modular PyTorch simulator for OTA-FL with noisy analog
    aggregation, local client training, server-side optimization, logging,
    result export, and plotting.
    \item Implementations of FedAvg-OTA, FedAvgM-OTA, AdaGrad-OTA, and
    Adam-OTA as server-side optimizers that consume the OTA-aggregated
    update.
    \item Alpha-stable tail-index ablations that test how OTA-FL behaves as
    the channel changes from AWGN at $\alpha=2$ to heavier-tailed
    interference at $\alpha<2$.
    \item A MAC versus no-MAC comparison under heavy-tailed interference,
    testing whether robust pre-processing complements adaptive server
    normalization.
    \item A lightweight one-dimensional theory sketch explaining why
    second-moment accumulation can induce automatic step-size shrinkage under
    impulsive channel noise.
\end{itemize}

\section{Thesis Organization}
Chapter 2 reviews FL, OTA aggregation, adaptive optimization, and the
relationship between this project and the reference paper. Chapter 3
describes the implemented system model and optimizer equations. Chapter 4
presents the datasets, experimental settings, generated plots, and numerical
results. Chapter 5 summarizes the outcome, limitations, and future work.
